\chapter{Accéder aux marchés~: le change}\label{ch:m2:getting-access-fx}

Le 15 janvier 2015, la Swiss National Bank a cessé de plafonner le franc, et le franc
a bondi (chapitre 18). Selon la BRI, ces mouvements ont secoué le secteur du prime brokerage,
industrie, les banques qui se portent garantes des transactions de leurs clients, et les prime brokers
ont réévalué leur activité~: ils ont conservé leurs plus gros clients, les grandes
sociétés de trading pour compte propre, et ont abandonné les agrégateurs de particuliers, les petits hedge funds et certains
traders haute fréquence, ont accru les exigences de fonds propres, durci les conditions d'admission et
relevé leurs commissions~; le volume FX traité par les prime brokers a reculé de 22~\% entre 2013 et 2016. Un nouveau fonds
dont la stratégie négocie un milliard de dollars par jour doit d'abord trouver quelqu'un qui se porte garant de
ses transactions, puis respecter les limites que cette personne fixe. Ce chapitre explique
comment~: les lignes de crédit du prime broker et les \omterm{def:m2:getting-access-fx:giveup}{lignes de give-up}, les limites et la passerelle qui les applique,
devenir fournisseur de liquidité, ainsi que les obligations d'information et les revues des deux côtés.

\section{Crédit du prime broker et lignes de give-up}

Un fonds ne peut pas négocier du FX avec des dizaines de dealers, dont chacun devrait évaluer son
crédit, signer la documentation et fixer une limite. Un prime broker FX (chapitre 14) s'en charge
une seule fois~: le fonds négocie avec tout dealer qui accepte le nom du prime broker, et chaque
transaction est \emph{remise en give-up} au prime broker, qui devient la contrepartie du dealer
et enregistre une transaction symétrique avec le fonds.

\begin{definition}[Ligne de give-up, avis de désignation]\label{def:m2:getting-access-fx:giveup}
Une \emph{ligne de give-up}\index{ligne de give-up} est un accord par lequel un prime broker accepte les
transactions qu'un client exécute auprès d'un dealer d'exécution donné, dans certaines limites. Un
\emph{avis de désignation}\index{avis de désignation} est l'avis du prime broker adressé à un
dealer d'exécution qui désigne le client comme partie autorisée à négocier pour son
compte, et fixe les types de transactions, les maturités maximales, les devises et les limites de cette
habilitation.
\end{definition}

Les documents types ont été publiés par le Foreign Exchange Committee à New York en
2005~: dans le Master FX Give-Up Agreement, le prime broker autorise chaque
partie désignée, par un avis établi selon un formulaire type, à négocier avec le dealer pour son compte~; la
habilitation est limitée aux types, maturités, devises et bureaux mentionnés dans l'avis et à
deux limites, et le prime broker n'est pas responsable d'une transaction qui porte une limite au-dessus
de son niveau, ou au-delà encore, sans son consentement
(\cref{fig:m2:getting-access-fx:giveup}). Le coût pour le client se compose des commissions du prime broker,
facturées par million négocié et pour le crédit accordé, ainsi que du droit du prime broker de réduire
les limites quand il le souhaite.

\begin{omfigure}
\centering
\begin{tikzpicture}[font=\small,
  box/.style={draw,thick,rounded corners=2pt,align=center,minimum height=1cm,minimum width=2.8cm,inner sep=3pt},
  arr/.style={-{Stealth[length=2.5mm]},thick}]
\node[box,draw=omDef,fill=omDef!8] (fund) at (0,0) {fonds\\(partie désignée)};
\node[box,draw=gray,fill=gray!8] (ed) at (5.4,1.4) {dealer d'exécution 1};
\node[box,draw=gray,fill=gray!8] (ed2) at (5.4,-1.4) {dealer d'exécution 2};
\node[box,draw=omThm,fill=omThm!8] (pb) at (10.8,0) {prime broker};
\draw[arr] (fund) -- node[above,sloped,font=\footnotesize] {exécute} (ed);
\draw[arr] (fund) -- node[below,sloped,font=\footnotesize] {exécute} (ed2);
\draw[arr] (ed) -- node[above,sloped,font=\footnotesize] {give-up} (pb);
\draw[arr] (ed2) -- node[below,sloped,font=\footnotesize] {give-up} (pb);
\draw[arr,dashed,omThm] (pb.north) |- node[pos=0.75,above,font=\footnotesize] {avis de désignation~: paires, maturités, NOP et limites de règlement} (ed.north);
\draw[arr,omThm] (pb.south) |- node[pos=0.75,below,font=\footnotesize] {transaction symétrique avec le fonds} ([yshift=-2.4cm]fund.south) -- (fund.south);
\end{tikzpicture}

\omcaption{Une \omterm{def:m2:getting-access-fx:giveup}{ligne de give-up}. Le fonds exécute auprès de tout dealer qui accepte le nom de son
prime broker~; chaque transaction est remise au prime broker, qui fait face au dealer et
enregistre une transaction symétrique avec le fonds. L'avis du prime broker indique à chaque dealer
ce que le fonds peut négocier pour son compte et dans quelles limites. Schéma de principe.}\label{fig:m2:getting-access-fx:giveup}
\end{omfigure}

\begin{definition}[Prime-of-prime]\label{def:m2:getting-access-fx:pop}
Le \emph{prime-of-prime}\index{prime-of-prime} est un dispositif dans lequel un client est
servi en prime brokerage par une banque ou une société non-dealer, qui est elle-même servie
en prime brokerage par une banque dealer FX.
\end{definition}

Après 2015, les clients délaissés par les banques se sont tournés vers des fournisseurs de \omterm{def:m2:getting-access-fx:pop}{prime-of-prime}, qui prennent
leur crédit et le regroupent en bundle dans leur propre relation avec une banque
(\cref{fig:m2:getting-access-fx:pop}). Il ajoute un niveau, avec ses coûts et ses limites, mais
c'est souvent le seul moyen d'accès pour une petite société. Le comité de New York a publié en 2014
modèles types de certification et bonnes pratiques pour ces dispositifs d'intermédiation, dans
lesquels l'intermédiaire certifie ses relations avec les prime brokers et les dealers.

\begin{omfigure}
\centering
\begin{tikzpicture}[font=\small,
  box/.style={draw,thick,rounded corners=2pt,align=center,minimum height=1cm,minimum width=2.6cm,inner sep=3pt},
  arr/.style={-{Stealth[length=2.5mm]},thick}]
\node[box,draw=omDef,fill=omDef!8] (c) at (0,0) {petit fonds ou\\agrégateur\\de particuliers};
\node[box,draw=omProp,fill=omProp!8] (pop) at (4.6,0) {prime-of-prime\\(hors dealers)};
\node[box,draw=omThm,fill=omThm!8] (pb) at (9.2,0) {banque dealer FX\\(prime broker)};
\node[box,draw=gray,fill=gray!8] (v) at (13.4,0) {plateformes\\et dealers};
\draw[arr] (c) -- node[above,font=\footnotesize] {crédit} (pop);
\draw[arr] (pop) -- node[above,font=\footnotesize] {crédit} (pb);
\draw[arr] (pb) -- node[above,font=\footnotesize] {nom} (v);
\end{tikzpicture}

\omcaption{\omterm{def:m2:getting-access-fx:pop}{Prime-of-prime}~: un client trop petit pour le prime brokerage d'une banque est
servi en prime brokerage par une société non-dealer, qui est elle-même servie par une banque dealer
et accède aux plateformes de négociation et aux dealers sous son nom. Schéma de principe.}\label{fig:m2:getting-access-fx:pop}
\end{omfigure}

\section{Limites~: position ouverte nette et règlement}

\begin{definition}[Limite de position ouverte nette, limite de règlement]\label{def:m2:getting-access-fx:limits}
Une \emph{limite de position ouverte nette}\index{limite de position ouverte nette} plafonne l'exposition
nette en devises du client auprès du prime broker, exprimée dans une seule devise pour l'ensemble
des devises. Une \emph{limite de règlement}\index{limite de règlement} plafonne, pour chaque date de valeur,
le montant que le prime broker devra régler au titre des transactions du client.
\end{definition}

Le Master Agreement définit les deux en valeurs en dollars~: la position ouverte nette agrège,
selon une méthode fixée dans l'annexe, la valeur nette en dollars de chaque devise sur l'ensemble
des transactions de la partie désignée~; le montant net quotidien de règlement d'une date de valeur est la somme
des valeurs en dollars des devises pour lesquelles un montant net est dû au prime broker ce
jour. La passerelle de ce chapitre retient comme position ouverte nette la somme des valeurs longues nettes en
dollars, qui égale la somme des valeurs courtes nettes lorsque chaque devise, le dollar
compris, est comptée, et le montant de règlement comme la somme des devises à
recevoir à la date de valeur.

\begin{example}[Deux transactions passant par la passerelle]\label{ex:m2:getting-access-fx:nop}
Le prime broker A autorise EURUSD et USDJPY jusqu'à sept jours, une position ouverte nette de 100
millions de dollars et un règlement de 500 millions de dollars par date de valeur. Le fonds achète 50 millions d'euros
à 1.10 au comptant~: il est long de 50 millions d'euros valant 55 millions de dollars et short de 55 millions de dollars, soit une
position ouverte nette de 55 millions et un montant de règlement de 55 millions. Il achète ensuite 20
millions de dollars contre yen~: la position courte en dollars se réduit à 35 millions et apparaît une position courte en yens de 20
millions de yens~; les longs restent ceux de 55 millions d'euros, et il en va de même du montant
de règlement. Un nouvel achat de 95 millions d'euros porterait la position à 159.5 millions
et il est rejeté~; une vente d'euros serait acceptée même si la limite avait été abaissée sous
la position, car elle la réduit.
\end{example}

Une limite ne vaut que par son application. La passerelle propre au fonds doit contrôler chaque ordre
avant son envoi, dans le parcours de l'ordre, au regard de ce qu'autorise l'avis de chaque prime broker~:
un dealer qui accepte une transaction au-delà de la limite peut voir le prime broker la refuser,
le fonds se retrouvant face au dealer directement, ou pas du tout. La construction de ce chapitre
est cette passerelle, écrite en C++20 et en Rust pour le parcours de l'ordre et en Python pour l'analyse
(voir la \cref{fig:m2:getting-access-fx:path}).

\begin{omfigure}
\begin{tikzpicture}
\begin{axis}[width=11cm,height=5cm,
  xlabel={ordres dans la journée},ylabel={NOP utilisée (\% de la limite)},
  tick label style={font=\small},label style={font=\small},
  xmin=0,xmax=240,ymin=0,ymax=105,grid=major,grid style={gray!25},
  legend columns=3,legend style={font=\footnotesize,at={(0.5,-0.3)},anchor=north,draw=none,fill=none,/tikz/every even column/.append style={column sep=8pt}},
  table/col sep=comma]
\addplot[omDef,thick] table[x=order,y=A]{figdata/markets-2/27-getting-access-fx/nop_path.csv};
\addplot[omThm,thick] table[x=order,y=B]{figdata/markets-2/27-getting-access-fx/nop_path.csv};
\addplot[omProp,thick] table[x=order,y=C]{figdata/markets-2/27-getting-access-fx/nop_path.csv};
\legend{A (frais 3),B (frais 5),C (frais 8)}
\end{axis}
\end{tikzpicture}

\omcaption{Une journée simulée de 240 ordres acheminés par la passerelle vers le prime broker le moins cher
dont les limites les autorisent~: la NOP de chacun, en pourcentage de sa limite. A
exécute d'abord~; à partir de l'ordre 107, le flux qui augmente la position déborde sur B, puis sur
C, dont la position termine la journée à 76~\% de sa limite. Aucun ordre n'est rejeté.
Illustratif~; données~: le tutoriel du chapitre, à graine fixée.}\label{fig:m2:getting-access-fx:path}
\end{omfigure}

\section{Devenir fournisseur de liquidité}

Un fonds qui affiche des cotations, plutôt que de seulement prendre les prix, doit aussi être accepté comme un
fournisseur de liquidité~: sur une plateforme de négociation, ce qui lui permet de diffuser ses prix dans le carnet d'ordres de la plateforme~; ou
directement auprès de ses clients, par l'intermédiaire d'une \omterm{def:m2:the-fx-market:sdp}{plateforme mono-dealer} ou d'un agrégateur.

\begin{definition}[Revue des fournisseurs de liquidité]\label{def:m2:getting-access-fx:review}
Une \emph{revue des fournisseurs de liquidité}\index{revue des fournisseurs de liquidité} est l'évaluation périodique,
par une plateforme ou un client, de la cotation d'un fournisseur de liquidité~: la fréquence à laquelle ses prix
sont touchés et exécutés, son \omterm{def:m2:fx-spot-microstructure:lastlook}{taux de rejet} et ses temps de rétention sous \omterm{def:m2:fx-spot-microstructure:lastlook}{last look}, ainsi que les
mark-outs des transactions réalisées avec lui.
\end{definition}

Les indicateurs sont ceux du chapitre 15, vus de l'autre côté. Un fournisseur de liquidité
qui rejette trop souvent, retient les demandes trop longtemps, ou dont les exécutions sont systématiquement
suivies de mouvements défavorables à ses clients peut être rétrogradé dans le routage des agrégateurs ou
retiré du flux d'une plateforme~; celui qui affiche des prix fermes et largement exécutés peut espérer davantage de flux.
Il lui faut son propre crédit pour cela~: chaque transaction qu'il conclut en tant que fournisseur de liquidité reste
donnée en give-up à son prime broker et compte dans les mêmes limites.

\section{Courtage, transparence et revues}

Le \omterm{def:m2:fx-spot-microstructure:code}{FX Global Code} demande aux intervenants de publier leurs modes de négociation, et en août
2021 le Global Foreign Exchange Committee a publié des fiches types de transparence
destinées aux fournisseurs de liquidité et aux plateformes de négociation, accompagnées d'un document d'orientation sur le \omterm{def:m2:fx-spot-microstructure:lastlook}{last look}.
La fiche d'un fournisseur de liquidité indique en un seul endroit s'il agit en principal ou
en agent, s'il partage les données issues des interactions avec ses clients, s'il pré-couvre, s'il
et comment il utilise le \omterm{def:m2:fx-spot-microstructure:lastlook}{last look} (de façon symétrique ou non, fenêtre maximale et minimale, et
s'il négocie pendant la fenêtre), et où trouver ses politiques d'agrégation,
de discrétion, d'horodatage, d'ordres stop-loss, d'exécutions partielles, de prix de référence, de mark-ups
et du partage interne d'informations confidentielles. Un fonds les lit avant de
choisir avec qui négocier, et doit rédiger les siennes s'il devient fournisseur de liquidité.

\begin{dated}{2026-09}{Documents d'accès au FX}\label{dat:m2:getting-access-fx:disclosure}
Master FX Give-Up Agreement~: publié par le Foreign Exchange Committee en 2005, avec un
formulaire type d'avis. \omterm{def:m2:fx-spot-microstructure:code}{FX Global Code}~: dernière mise à jour en décembre 2024, avec une fiche de transparence des fournisseurs de liquidité de décembre
2024~; fiches publiées pour la première fois en août
2021.
\end{dated}

\section{Tutoriel~: une passerelle de limites pré-négociation}\label{tut:m2:getting-access-fx:gate}

\begin{tutorial}
\textbf{Objectif.} Implémenter un contrôle pré-négociation qui fait respecter, pour chaque prime broker, les
paires de devises et les échéances autorisées, les \omterm{def:m2:getting-access-fx:limits}{limites de position ouverte nette} et de règlement, acheminer chaque ordre vers
le prime broker le moins cher qui l'accepte, et planifier la répartition du flux d'une journée.
\textbf{État final~:} les \cref{fig:m2:getting-access-fx:path,fig:m2:getting-access-fx:plans},
l'\cref{ex:m2:getting-access-fx:nop}, ainsi que les chiffres du problème du week-end.
\begin{tutsteps}
\item \textbf{La passerelle}, en C++20~: le contrôle, qui laisse le registre inchangé, et la
mise à jour une fois la transaction conclue.
\omcode{code/firm/pblimits/cpp/firm_pblimits.hpp}{46}{75}{Le registre du prime broker et son contrôle pré-négociation, en C++20.}\label{lst:m2:getting-access-fx:gate}
\item \textbf{Le routeur} et le fichier des limites, en Python.
\omcode{code/firm/pblimits/firm_pblimits.py}{85}{102}{Routage vers le prime broker le moins cher disposant encore de marge, et lecture des limites.}\label{lst:m2:getting-access-fx:route}
\item \textbf{Exécuter} \texttt{access\_demo.plan()}, \texttt{access\_demo.brute\_force()},
\texttt{access\_demo.day\_of\_orders()} et \texttt{fig\_access.py}~; construire et tester les
passerelles C++ et Rust.
\end{tutsteps}
\textbf{Pour aller plus loin.} Ajouter une taille maximale d'ordre par paire à chaque avis~; puis
faire en sorte que le routeur privilégie, parmi les prime brokers qui acceptent un ordre, celui dont
le coût de portage overnight est le plus bas lorsque l'ordre augmente la position, et mesurer le coût
économisé.
\end{tutorial}

\section{Construction~: la passerelle de limites des prime brokers}\label{bld:m2:getting-access-fx:pblimits}

\begin{build}
\bfield{Rôle} Chaque ordre que la société miniature envoie en FX passe par cette passerelle avant
de quitter le fonds~: il doit tenir dans l'avis d'un prime broker, et il est confié en give-up au
moins cher qui l'accepte.
\bfield{Interface} C++20 \texttt{firm::pblimits::Book} avec \texttt{check(order,
usd\_per)}, \texttt{apply}, \texttt{nop()}, \texttt{settlement(day)}, et
\texttt{route(books, order, usd\_per)}~; la même chose en Rust~; en Python \texttt{Book},
\texttt{check}, \texttt{apply}, \texttt{route} et \texttt{read\_limits(path)} pour le
fichier de limites partagé \texttt{code/firm/pblimits/limits.csv}.
\bfield{Règles} Valeurs en dollars aux taux donnés~; position ouverte nette égale à la somme des longs nets~;
règlement égal à la somme des encaissements nets par date de valeur~; rejeter si une paire ou une échéance n'est pas
autorisée, ou si une limite est dépassée ou dépassée davantage~; acheminer par frais croissants.
\bfield{Tests d'acceptation} Les mêmes dans les trois langages~: la position ouverte nette et
le règlement de deux transactions~; rejet au-dessus de la limite et acceptation d'une transaction réductrice
après une baisse de limite~; la \omterm{def:m2:getting-access-fx:limits}{limite de règlement} par date de valeur~; rejets de paires et d'échéances~; routage
vers le courtier le moins cher disposant de marge.
\bfield{Pour aller plus loin} Limites par dealer d'exécution aussi bien que par prime broker~; un
registre sans verrou pour les fils concurrents d'ordres (livre~13)~; transferts de
positions de fin de journée entre prime brokers~; limites sur les options par leur delta et leur prime.
\end{build}

\begin{omsources}
\item Foreign Exchange Committee, Master FX Give-Up Agreement (2005)~; FXC et FMLG,
notes de pratique de marché de mai 2013 et de juillet 2014 sur les avis des prime brokers et
dispositifs d'intermédiation.
\item M.~Moore, A.~Schrimpf et V.~Sushko, ``Downsized FX markets: causes and
implications'', BIS Quarterly Review, décembre 2016.
\item Global Foreign Exchange Committee, communiqué de presse du 18 août 2021 et
fiches de présentation des informations~; une fiche publiée d'un fournisseur de liquidité.
\end{omsources}

\section{Exercices}

\begin{exercise}[$\star$]\label{exo:m2:getting-access-fx:1}
Le portefeuille d'un fonds est long en EUR pour une valeur de 40~millions de dollars, court de 25~millions de dollars et court en JPY pour une valeur de
15~millions de dollars. Quelle est sa position ouverte nette, et pourquoi le total des longs est-il égal à celui des courts~?
\end{exercise}

\begin{exercise}[$\star$]\label{exo:m2:getting-access-fx:2}
Qu'indique un \omterm{def:m2:getting-access-fx:giveup}{avis de désignation} à un dealer d'exécution~?
\end{exercise}

\begin{exercise}[$\star$]\label{exo:m2:getting-access-fx:3}
Pourquoi les prime brokers ont-ils délaissé leurs clients les plus petits après 2015, et vers où ces clients se sont-ils reportés~?
\end{exercise}

\begin{exercise}[$\star\star$]\label{exo:m2:getting-access-fx:4}
Pour une même date de valeur, le fonds achète 100~millions d'euros et en vend 60~millions d'euros contre
dollars au taux de 1.10. Quel est le montant de règlement-livraison, et quel serait-il si les deux transactions
étaient passées auprès de prime brokers différents~?
\end{exercise}

\begin{exercise}[$\star\star$]\label{exo:m2:getting-access-fx:5}
Pourquoi le filtre de contrôle accepte-t-il une transaction qui réduit une position déjà au-dessus de sa limite~?
\end{exercise}

\begin{exercise}[$\star\star$]\label{exo:m2:getting-access-fx:6}
Quels chiffres une plateforme de négociation examinerait-elle pour procéder à la revue du fonds en tant que fournisseur de liquidité~?
\end{exercise}

\begin{exercise}[$\star\star\star$]\label{exo:m2:getting-access-fx:7}
\emph{Programmation.} Ramener la \omterm{def:m2:getting-access-fx:limits}{limite de position ouverte nette} du prime broker A à 50~millions de dollars et relancer la
journée d'ordres. Quel flux supplémentaire part vers C, et où finit la
position de la journée~?
\end{exercise}

\begin{exercise}[$\star\star\star$]\label{exo:m2:getting-access-fx:8}
\emph{Repérer l'erreur.} ``Notre stratégie ne négocie que des devises majeures liquides~; un seul prime broker
suffit.'' La corriger.
\end{exercise}

\section{Problème~: Trois prime brokers}

\begin{problem}[{Problème du week-end --- garder chaque limite au moindre coût}]\label{pb:m2:getting-access-fx:1}
La journée d'un fonds, en millions de dollars~: 1\,000 de flux sur EURUSD et USDJPY, 150 sur GBPUSD et 50
sur USDMXN, et une position ouverte nette de nuit de 300, dont 30 en pesos. Trois
avis des prime brokers, extraits de \texttt{limits.csv}~: A facture 3 dollars par million négocié,
autorise EURUSD et USDJPY jusqu'à 7 jours, une position ouverte nette de 100 et un règlement-livraison de 500~;
B facture 5, ajoute GBPUSD, accorde une durée maximale de 30 jours, 150 et 600~; C facture 8, ajoute USDMXN, accorde
une durée maximale d'un an, 300 et 1\,000. Les positions de nuit coûtent 30, 20 et 50 dollars par million et par jour
chez A, B et C. Pour la planification, prendre le montant de règlement-livraison chez un prime broker égal à
l'ensemble du flux confié en give-up à celui-ci, sans compensation.

\medskip\noindent\textbf{Partie I --- Les avis.}
\begin{enumerate}
\item Quels prime brokers peuvent prendre chaque partie du flux~?
\item Calculer la position ouverte nette après l'achat de 50~millions d'euros à 1.10, puis de
20~millions de dollars contre le yen.
\item Chez B, après l'achat de 500~millions d'euros au comptant, le fonds peut-il en acheter 50~millions d'euros de plus pour
la même date de valeur~? Et pour le jour suivant~?
\item La limite de A est ramenée à 50~millions alors que le fonds est long de 99~millions~: quels ordres
passent encore~?
\item Pourquoi l'avis limite-t-il les échéances~?
\end{enumerate}

\medskip\noindent\textbf{Partie II --- Le plan.}
\begin{enumerate}[resume]
\item Répartir le flux au moindre coût dans les \omterm{def:m2:getting-access-fx:limits}{limites de règlement}.
\item Répartir la position de nuit au moindre coût dans les \omterm{def:m2:getting-access-fx:limits}{limites de position ouverte nette}.
\item Quel est le coût du plan par jour, et par année de 250~jours~?
\item Combien coûterait de tout traiter chez C, le seul prime broker qui accepte toutes les
devises~?
\item Vérifier le plan par rapport à une recherche sur toutes les répartitions par pas de 10~millions.
\end{enumerate}

\medskip\noindent\textbf{Partie III --- La journée.}
\begin{enumerate}[resume]
\item Acheminer une journée simulée de 240~ordres vers le prime broker le moins cher disposant de marge~: quand la
position de A atteint-elle pour la première fois sa limite~?
\item Combien d'ordres sont-ils rejetés~?
\item Où se termine la position de la journée, et pourquoi chez le prime broker le plus cher~?
\item Pourquoi les montants de règlement-livraison du routeur peuvent-ils être très inférieurs au flux qu'il envoie à un prime broker~?
\item Comment faudrait-il rapprocher le routeur du plan~?
\end{enumerate}

\medskip\noindent\textbf{Partie IV --- Jugement.}
\begin{enumerate}[resume]
\item Que se passe-t-il pour le fonds si B réduit de moitié ses limites après un choc~?
\item Pourquoi un dispositif de \omterm{def:m2:getting-access-fx:pop}{prime-of-prime} est-il une alternative, et qu'apporte-t-il~?
\item Que faudrait-il publier, et consulter, avant de devenir fournisseur de liquidité~?
\item Énoncer le \emph{résultat nommé}~: la répartition du flux et de la position qui garde chaque
limite au moindre coût, et ce coût.
\item En une phrase~: de quoi un fonds a-t-il besoin avant sa première transaction FX~?
\end{enumerate}
\end{problem}

\begin{omfigure}
\begin{tikzpicture}
\begin{axis}[width=11cm,height=4.8cm,ybar stacked,bar width=26pt,
  ylabel={coût par jour (milliers USD)},
  tick label style={font=\small},label style={font=\small},
  xtick=data,xticklabels={tout chez C,{frais\\optimisés},{NOP\\optimisée},{les deux\\optimisés}},xticklabel style={align=center},
  ymin=0,ymax=27,grid=major,grid style={gray!25},enlarge x limits=0.18,
  legend columns=2,legend style={font=\footnotesize,at={(0.5,-0.38)},anchor=north,draw=none,fill=none,/tikz/every even column/.append style={column sep=8pt}},
  table/col sep=comma]
\addplot[fill=omDef!60,draw=omDef] table[x=k,y=fees]{figdata/markets-2/27-getting-access-fx/plans.csv};
\addplot[fill=omThm!50,draw=omThm] table[x=k,y=nop]{figdata/markets-2/27-getting-access-fx/plans.csv};
\legend{frais sur le flux,frais sur les positions de nuit}
\end{axis}
\end{tikzpicture}

\omcaption{Coût quotidien du fonds du problème du week-end selon quatre plans~: tout chez C~;
le flux réparti au moindre coût de frais, la position restant chez C~; la position répartie au
moindre coût de charges, le flux restant chez C~; et les deux répartis. Le meilleur plan coûte 13\,800~USD par
jour contre 24\,600~USD. Illustratif~; données~: le tutoriel du chapitre.}\label{fig:m2:getting-access-fx:plans}
\end{omfigure}

\section{Questions d'entretien}

\begin{interviewq}[$\star$ \iqroles{trader, développeur}]\label{iq:m2:getting-access-fx:1}
Que fait un prime broker FX pour un fonds, et que signifie un give-up~?
\end{interviewq}

\begin{interviewq}[$\star$ \iqroles{risque, développeur}]\label{iq:m2:getting-access-fx:2}
Quelles sont la \omterm{def:m2:getting-access-fx:limits}{limite de position ouverte nette} et la \omterm{def:m2:getting-access-fx:limits}{limite de règlement}, et comment les calculeriez-vous~?
\end{interviewq}

\begin{interviewq}[$\star\star$ \iqroles{développeur}]\label{iq:m2:getting-access-fx:3}
À quel endroit du cheminement de l'ordre le contrôle de limite doit-il se placer, et que doit-il faire quand il
ne peut pas décider à temps~?
\end{interviewq}

\begin{interviewq}[$\star\star$ \iqroles{trader, risque}]\label{iq:m2:getting-access-fx:4}
Comment répartiriez-vous le flux d'un fonds entre plusieurs prime brokers~?
\end{interviewq}

\begin{interviewq}[$\star\star$ \iqroles{trader}]\label{iq:m2:getting-access-fx:5}
Que doit demander un fonds à un fournisseur de liquidité avant de négocier avec lui~?
\end{interviewq}

\begin{interviewq}[$\star\star\star$ \iqroles{développeur}]\label{iq:m2:getting-access-fx:6}
Concevez une porte de limite qui contrôle 100\,000 ordres par seconde répartis sur dix prime brokers, avec
une latence maximale de cinq microsecondes.
\end{interviewq}
